\documentclass[a4paper,11pt,spanish,sans]{exam}
\usepackage[spanish]{babel}
%\usepackage[utf8]{inputenc}
\usepackage{multicol}
%\usepackage[latin1]{inputenc}
\usepackage{fontspec}%la posta para las tildes con lualatex
\usepackage[margin=0.4in]{geometry}
\usepackage{amsmath,amssymb}
\usepackage{multicol}
\usepackage{natbib}
\usepackage{graphicx}
\usepackage{hyperref}
\usepackage{epstopdf}
\usepackage{capt-of}
\usepackage[usenames]{color}
\usepackage{pgf,tikz}
\usepackage{mathrsfs}
\usetikzlibrary{arrows}
\usetikzlibrary{babel}
\usepackage{pst-fractal}
\usetikzlibrary{decorations.markings}
\usetikzlibrary{shapes.geometric}
\usetikzlibrary{shapes,snakes}
\usepackage{tkz-euclide}
\usetkzobj{all}
\usepackage{fullpage} 
\usepackage{parskip}
\usepackage{caption}
\usepackage{subcaption}
%los de aca abajo capaz no los uso
\newcommand{\class}{Matemática: Guía de Números Complejos}
\newcommand{\term}{3° Trimestre 2015}
\newcommand{\examnum}{}
\newcommand{\examprof}{Alexis Gomel}
\newcommand{\examdate}{15/7/2015}
\newcommand{\timelimit}{60 Minutes}%no lo uso
\newcommand{\webpdf}{https://drive.google.com/file/d/0B2MOYme4kZd-OVVPdzF4VGoya1E/view?usp=sharing}%no lo uso
\newcommand{\Ts}{\rule{0pt}{2.6ex}}       % Top strut
\newcommand{\Bs}{\rule[-1.2ex]{0pt}{0pt}} % Bottom strut

%el header de las hojas.
\pagestyle{head}
\firstpageheader{}{}{}
\runningheader{\class}{$\qquad$ \examnum\ - Pagina \thepage\ de \numpages}{\examdate}
\runningheadrule


\begin{document}


\noindent
\begin{tabular*}{\textwidth}{l @{\extracolsep{\fill}} r @{\extracolsep{6pt}} l}
\textbf{\class} & \textbf{Profesor: \examprof}\\

\textbf{PDF: \href{\webpdf}{\textcolor{blue}{http://tinyurl.com/complejs}}} %& Teaching Assistant & \makebox[2in]{\hrulefill}
\end{tabular*}\\
\rule[2ex]{\textwidth}{2pt}

\[e^{i\pi}+1=0\]

Ahora que tengo su atención...

{{\small Esta guía es simplemente una guía. \textbf{\underline{NO} reemplaza ni incluye todo el material que se da en clase}}}.

\begin{center}
\section*{Guía de Números Complejos}
\end{center}
%vimos que enteros.. reales.. etc... buenos. ahora vamos a ver otros mas generales... los complejos.
%ver definicion de mate 4

\section*{Numero Imaginario:} 

\begin{minipage}{0.5\textwidth}
Se define al numero Imaginario como un numero "$i$" que cumple que 
\[i^2=-1\]

O también: 
\[i=\sqrt{-1}\] 
\end{minipage}
\begin{minipage}{0.5\textwidth}
	
	\centering
	\includegraphics[width= 1\linewidth]{imwtf.jpeg}
	%\caption{Set joke}
	\label{fig:erise}
	
\end{minipage}

\section*{Números Complejos:}

$z=a+ib=Re(z)+i.Im(z) \qquad $
$\begin{cases}
$a\equiv Re(z)$ \quad \text{'Parte Real de z'}\\
$b\equiv Im(z)$ \quad \text{'Parte imaginaria de z'}\\
\end{cases}$

ejemplo: $z=2+3i$


\subsection*{Representación Gráfica:}

Cada numero complejo se puede representar como un punto en el plano. Donde el eje x denota el conjunto de los números Reales , y el eje y el conjunto de los números Imaginarios.

$Re(z) \rightarrow  x \qquad Im(z) \rightarrow y $
\begin{figure}[!htbp]
	\begin{center}
		\begin{tikzpicture}
		\begin{scope}[thick,font=\scriptsize]
		
		%\draw[step=1cm,gray,dashed, very thin] (-3.9,-3.9) grid (3.9,3.9);
		
		\draw[->] (-4,0) -- (4,0) node [below left] {$Re\{z\}$};
		\draw[->] (0,-4) -- (0,4) node [above right] {$Im\{z\}$};
		
		\foreach \n in {-3,...,-1,1,2,...,3}{%
			\draw (\n,-3pt) -- (\n,3pt)   node [above] {$\n$};
			\draw (-3pt,\n) -- (3pt,\n)   node [right] {$\n i$};}
		
		\draw [thick, color=red] (0,0) -- (2,3);
		\draw [color=blue, fill=blue] (2,3) circle(0.05);
		\node [color=black] at (3,3.2) {$2+3i=(2,3)$};
		\end{scope}
		\end{tikzpicture}
	\end{center}
	\caption{Represantacion grafica del numero \textit{\textbf{2+3i}} en el Plano Complejo (diagrama de Argand).}
	
\end{figure}

\section{Operaciones:}

Propiedades:
Al igual que con los Reales, los Complejos cumplen la mayoría de las propiedades que solemos usar con los números reales respecto de la suma y el producto:

Suma y producto: Neutro, Inverso, conmutatividad, asociatividad.

\begin{itemize}
	\item Suma:
	\[z_1=a+ib \qquad z_2=c+id \qquad \Rightarrow z_1+z_2=(a+c) + i(b+d)\]
	
	ejemplo: $ \begin{cases}
	z_1=1+2i\\ 
	z_2=2+3i\\
	\end{cases}$  $\qquad \Rightarrow   z_1+z_2=3+5i$
	
	\item Multiplicacion:
	\[z_1=a+ib \qquad z_2=c+id \qquad \Rightarrow z_1.z_2=(a+ib).(c+id)= (ac-bd)  +  i(ad+bc)\]
	
	ejemplo: $ \begin{cases}
	z_1=1+2i\\ 
	z_2=2+3i\\
	\end{cases}$  $\qquad \Rightarrow z_1.z_2=2-6+i(3+4)=-4+7i$
	
	\item División:
	\[z_1=a+ib \qquad z_2=c+id \qquad \Rightarrow \dfrac{z_1}{z_2}=\dfrac{a+ib}{c+id}=\dfrac{a+ib}{c+id}\cdot\dfrac{c-id}{c-id} =\dfrac{(a+ib)(c-id)}{c^2+d^2}\]
	
	\item Conjugación:
	
	Sea $z=a+ib$. 
	
	Se le llama $z^*$ al conjugado de $z$: \[z^*=a - ib\]
	
	Ejemplo: $z_1=1+2i$:  $\qquad \Rightarrow z_1^*=1-2i$
	
	 
\end{itemize}


\subsection*{Modulo y fase de un numero complejo:}

Se le llama 'Modulo', 'Norma ' o distancia de un numero complejo, a la distancia al origen de coordenadas del numero $z$, y se escribe como:
\[|z|=\sqrt{a^2+b^2}=\sqrt{z.z^*}\]

Se le llama 'fase' o 'argumento' de un numero complejo, al angulo que forma el vector del numero complejo con el eje de los números Reales.
\[\alpha=\tan^{-1} \left( \frac{b}{a} \right) \]



\begin{figure}
	\begin{center}
		\begin{tikzpicture}
		\begin{scope}[thick,font=\scriptsize]
		
		%\draw[step=1cm,gray,dashed, very thin] (-3.9,-3.9) grid (3.9,3.9);
		
		\draw[->] (-4,0) -- (4,0) node [below left] {$Re\{z\}$};
		\draw[->] (0,-4) -- (0,4) node [above right] {$Im\{z\}$};
		
		\foreach \n in {-3,...,-1,1,2,...,3}{%
			\draw (\n,-3pt) -- (\n,3pt)   node [above] {$\n$};
			\draw (-3pt,\n) -- (3pt,\n)   node [right] {$\n i$};}
		
		\draw [thick, color=red] (0,0) -- (-2,2);
		\draw [color=blue, fill=blue] (-2,2) circle(0.05);
		\draw [color=black, dotted , opacity=.2] (0,0) circle(2.828);
		\draw [dotted, color=black, opacity=0.2] (-2,0) -- (-2,2);
		
		\coordinate (d) at (-2,2);
		
		\draw[color=red]
		(0.5,0) coordinate (a) node[right] {}
		-- (0,0) coordinate (b) node[left] {}
		-- (-0.5,0.5) coordinate (c) node[above right] {};
		\tkzLabelSegment[right=2pt](b,d){$|z|$}
		\tkzMarkAngle[fill=orange,size=0.6,opacity=0.1](a,b,c);
		\tkzLabelAngle[pos = 0.35](a,b,c){$\hat{\alpha}$}
		
		\node [color=black] at (-2.3,2.3) {$ -2+2i=\sqrt{8}(\cos(\alpha)+i\sin(\alpha))$};
		\end{scope}
		\end{tikzpicture}
	\end{center}
\end{figure}

 \subsection*{Representaciones:}
\begin{itemize}
	\item Binomica o Cartesiana: $z=(a, b)$
	\item Algebraica: $z=a+ib$
	\item Polar:
	$z=\sqrt{a^2+b^2}.(\cos(\alpha)+i.\sin(\alpha))=|z|.(\cos(\alpha)+i.\sin(\alpha))$ 
	
	Donde $\alpha$ es el angulo de hace el numero complejo con el eje de los números Reales ($\alpha=\tan^{-1}\left( \frac{b}{a}\right) $).
	
\end{itemize}

Ejemplo: $1+i=(1,1)=\sqrt{2}(\cos(45\degrees)+i.\sin(45\degrees))$

\section*{Ejercicios}

TP: Los que necesiten nota, pueden entregar esta guía de ejercicios resuelta como trabajo practico. Para aprobar el trabajo practico hace falta resolver (justificados) almenos el 70 porciento de cada sección. 

Los ejercicios con $\bigstar $ son obligatorios, menos el que tiene $\bigstar \bigstar$.

Si se traban con algún ejercicio, pasen al siguiente, y vuelvan al ejercicio difícil mas tarde.

\section{Completar:}
\begin{center}
	%\caption{Completar}
	\label{completar}
	\begin{tabular}{|c|c|c|c|c|c|c|c|c|}
		\hline
		$i^2$  & $i^3$  &  $i^4$ &  $i^5$ &  $i^{16}$ & $i^{10}$ & $i^{-1}$  & $i^{4.n} \quad n\in \mathbb{N}$ & $i^{4.n+1} \quad n\in \mathbb{N}$  \Ts \Bs     \\ \hline
		 &  &  &  &  &  &  & &     \Ts \Bs     \\ \hline  
	\end{tabular}
	%\end{table*}
\end{center}

\section{Graficos:}
Representar los siguientes números en su forma  cartesiana y  polar. Graficarlos en el plano complejo.
\begin{multicols}{2}
\begin{enumerate}%cambiar funcion %ojo con cuales son asintotas oblicuas
\item $z=2-i$
\item $z=1+i$
\item $z=-2+3i$
\item $z=-1-i$

\columnbreak

\item $z=2(\cos(0)+i. \sin(0))$
\item $z=2(\cos(\pi /3)+i. \sin(\pi /3))$
\item $z=2(\cos(7\pi /4)+i. \sin(7\pi /4))$
\item $z=\frac{1}{2}(\cos(\pi /2)+i. \sin(\pi/2))$

\end{enumerate}
\end{multicols}

\section{Calcular: }

\begin{multicols}{2}
	\begin{enumerate}
		\item $(2+5i) + (-3+4i) + (2-8i)$
		\item $(1/2+i1/3)+(-1/3+i1/2)$
		\item $(2/3-i1/2)-(-3/2+i1/3)$
		\item $(5+5i).(-3+5i)$
		\item $(5-3i).(-3+5i)$
		\columnbreak
		\item $\dfrac{3-2i}{2+i}$
		\item $\dfrac{1+i}{3-i}$
		\item $\dfrac{(-1+i)+(2-4i)}{-2+3i}$
		\item $\dfrac{(5-2i).(1-i)}{-1-i}$
	\end{enumerate}
\end{multicols}

\section{Graficar y escribir las cuadraticas en todas sus formas:}
(Les sirve de repaso para la integradora)

\begin{minipage}{0.45\textwidth}
\begin{enumerate}
	\begin{multicols}{2}
	\item $x^2+1$
	\item $x^2+4x+5$
	\columnbreak
	\item $(x-1)^2+1$
	\item $-2x^2+8x-11$
\end{multicols}
\end{enumerate}
\end{minipage}
\begin{minipage}{0.65\textwidth}
	\centering
	\includegraphics[width= 1\linewidth]{pics.jpg}
	%\caption{Set joke}
	\label{fig:erise}
\end{minipage}

\section{Resolver}
	
	\begin{enumerate}
		\begin{minipage}{0.7\textwidth}
		\begin{multicols}{2}
		\item $(2+3i)^2$
		\item $(2-i)^2$
		\item $(-3+2i)^2$
		\item $(-2-3i)^3$
		\item $(-3-2i)^3$
		\item $\dfrac{(1+i)^2}{1-i}$
		
		\columnbreak
		
		
		\item $\dfrac{(2-3i)^2}{2+i}$
		\item $\dfrac{3-i}{(1+i)^2}$
		\item $\dfrac{(5+i)i^2}{(-1-i)^2}$
		\item $\dfrac{\sqrt{6}-\sqrt{10}i}{\sqrt{3}+\sqrt{2}i}$
		\end{multicols}
	\end{minipage}
		\begin{minipage}{0.25\textwidth}
			
			\centering
			\includegraphics[width= \linewidth]{xkcd.png}
			%\caption{Set joke}
			\label{fig:erise}
			
		\end{minipage}
		
	$\bigstar$	Dados $\quad z_1=2-1$ ;  $\quad z_2=-2+3i$ ;  $\quad z_3=1-i$.
	
		\begin{multicols}{2}
			\item $(z_1+z_2)^2$
			\item $\dfrac{z_1}{z_2.z_3}$
			\item $|z_1|^2+\left(  \dfrac{z_2}{z_3} \right) ^{-1}$
			
			\columnbreak
			
			\item $z_1+z_1^*$
			\item $\dfrac{z_2^2+|z_{3}|^2}{z_1^*}$
			\item $i^{43}.z_1+\dfrac{i^{37}}{z_2}-i^{122}.z_3^*$
		\end{multicols}
	\end{enumerate}
	
		
	

\section{ Ecuaciones:}

$\bigstar$ Calcular los valores de x e y para que es cumplan las siguientes igualdades
\begin{enumerate}

\begin{multicols}{2}
	\item $(2x; y+2)=(4; -1)$
	\item $3x-1+(1-y)i=(2;3)$
	\item $(-\frac{1}{2}x+3;-y+\frac{1}{4})=(0; 1)$
	\item $(2x-5)i-4y+1=3-i$

	
	\columnbreak
	
	\centering
	\includegraphics[width= 0.6\linewidth]{bereal.jpg}
	%\caption{Set joke}
	\label{fig:erise}
	
\end{multicols}

Hallar el valor de z:
	\item $i^{2}(5-3z)=z.i^{16}+4i$
	\item $-z(4+i.z)+1=(4+2i)z+i^4$
\end{enumerate}

\rule[2ex]{\textwidth}{2pt}

Observación: Para saber si hiciste bien un ejercicio, reemplaza por tu valor de x.

\section*{Extra y curiosidades}

\begin{itemize}

\item Una de las cosas mas llamativas y conocidas relacionadas con los numeros complejos es el conjunto o \href{https://www.youtube.com/watch?v=ZDU40eUcTj0}{\textcolor{blue}{Fractal de Mandelbrot}}. 
El cual fue 'descubierto' o 'inventado' por Benoit  mandelbrot en 1980 .

Un fractal es una elemento que repite alguna propiedad o forma a medida que uno hace zoom sobre el mismo. Se dice que un fractal es autosemejante, ya que cada parte de si mismo, se asemeja a si mismo.

\begin{minipage}{0.5\textwidth}
	\centering
	\includegraphics[width= 0.7\linewidth]{serptriang.png}
	%\caption{Cuadro de Escher}
	\label{fig:univerise}
\end{minipage}
\begin{minipage}{0.5\textwidth}
El mismo se consigue iterando para algun valor complejo 'c' la siguiente sucesion por recursion: 
\[\begin{cases}
z_0=0\\
z_{n+1}=z_{n}^2+c
\end{cases}
\]

Y separando a los puntos 'c' para los cuales la sucesion tiende a un valor determinado (converge), o sigue variando indefinidamente hacia el infinito (diverge).
Para los valores divergentes se le asigna un color a cada punto según que tan rápido crece su modulo .
\end{minipage}

Cerca de un conjunto de  valores, en el limite entre los dos casos, existe un conjunto de valores para los cuales la sucesión converge, pero cada punto o hace de manera distinta y a velocidades distintas, sin importar que tan cerca estén entre si.

Sin embargo lo hacen de tal manera que al graficarlos ciertas propiedades se repiten a medida que uno se 'acerca' (hace zoom) a un punto, asi indefinidamente.

Por ejemplo, la forma del escarabajo del mandelbrot se repite infinitas veces cerca de la frontera del escarabajo.

 

	
\begin{minipage}{0.5\textwidth}

\centering
\includegraphics[width= 0.7\linewidth]{mandel00.jpg}
%\caption{Cuadro de Escher}
\label{fig:univerise}

\end{minipage}
\begin{minipage}{0.5\textwidth}

\centering
\includegraphics[width= 0.7\linewidth]{mandel1.jpg}
%\caption{Set joke}
\label{fig:erise}

\end{minipage}
\begin{minipage}{0.5\textwidth}
	
	\centering
	\includegraphics[width= 0.7\linewidth]{mandel2.jpg}
	%\caption{Set joke}
	\label{fig:erise}
	
	\end{minipage}		
\begin{minipage}{0.5\textwidth}
	
	\centering
	\includegraphics[width= 0.72\linewidth]{Fractal_Broccoli.jpg}\\
	Romanesco Broccoli, una de las tantas cosas autosemejantes de la naturaleza.
	\label{fig:broccoli}
	
\end{minipage}	
				
					
El estudio del fractal de mandelbrot no solo nos da una idea de lo antintuitiva que puede ser la matemática aveces, sino que también tiene utilidades.
Por ejemplo, los gráficos de \href{https://www.youtube.com/watch?v=s65DSz78jW4}{\textcolor{blue}{videojuegos y muchas peliculas}} y muchas peliculas son generados a partir de fractales. Y muchas cosas en la naturaleza tiene propiedades parecidas a los fractales .

	
\end{itemize}


\newpage

	\begin{center}
		\section*{Axiomas de $\mathbb{R}$}
	\end{center}
	
	\begin{itemize}
		\subsection*{Suma:}
		
		\item S1. Conmutatividad: $a+b=b+a$ 
		\item S2. Asociatividad: $(a+b)+c=a+(b+c)=a+b+c$ 
		\item S3. Existe el Neutro de Suma : $\exists \: 0/ \quad a + 0=a$ 
		\item S4. Existe el inverso aditivo: $\forall a, \: \: \exists \: (-a) / \quad a + (-a)=0$ 
		
		
		\subsection*{Producto:}
		\item P1. Conmutatividad: $a.b=b.a$ 
		\item P2. Asociatividad: $(a.b).c=a.(b.c)=a.b.c$ 
		\item P3. Existe el Neutro del Producto: $\exists \: 1/ \quad a \cdot 1=a$  
		\item P4. Existe el inverso del Producto: $\forall a\neq \: 0, \: \: \exists \:  a^{-1} / \quad a.a^{-1}=1$   \\
		
		\item SP. Distributiva: $a.(b+c)=a.b+a.c$
		
		\subsection*{Axiomas de Orden:}
		\item O1. Tricotomia: Solo una de las siguientes posibilidades es verdadera: $a<b$, $a=b$, $a>b$
		\item O2. Transitividad: Si $a<b$ ; $b<c$ $\Rightarrow a<c$ (idem para el $=$ y el $>$)
		\item O3. Monotonía de Suma respecto del Orden: $a<b \Rightarrow a+c<b+c$ (idem para el $=$ y el $>$)
		\item O4. Monotonía del Producto respecto del Orden: $a<b \Rightarrow a.c<b.c$ (idem para el $=$ y el $>$)
		
		\subsection*{Axioma de Completitud:}
		\item C1. Sea $A \subset \mathbb{R} $, $A\neq \emptyset$, acotado superiormente $\Rightarrow \exists S=sup(A)$ (S = Supremo de A).
	\end{itemize}

\end{document}